\providecommand{\RAE}{\mathbb E}
\providecommand{\RAP}{\mathbb P}
\providecommand{\RAone}{\mathbf1}
\providecommand{\RAKL}{\operatorname{KL}}

\section{A finite direct assay for the mean boundary escape rate}
\label{sec:rate_assay}
The required decision statistic is an arithmetic average of actual local
operating rates. For a menu setting $u$ and target configuration $b$, let
\[
 \lambda_i^u=Q^u(b,b^{(i)}),\qquad
 \mu_u=\frac1N\sum_{i=1}^N\lambda_i^u,
 \qquad \Lambda_u=N\mu_u.
\]
Suppose a common certified envelope supplies
\begin{equation}
 0<\lambda_-\le\lambda_i^u\le\lambda_+<\infty,
 \qquad \kappa=\lambda_+/\lambda_-.
 \label{ra:bounds}
\end{equation}
For the preserved binary Gibbs kinetics with $b_i=1$,
$\lambda_i^u=r_i^u(b_{\partial i})[1-q_i^u(b_{\partial i})]$.
Thus $r_i^u\ge r_-$ and $q_i^u\le\bar q<1$ give the usable fast-direction
floor $\lambda_-=r_-(1-\bar q)$, and $r_i^u\le r_+$ gives
$\lambda_+=r_+$. These are bounds on the requested operating environments,
not rates transported from balancing fields.
Random-context averaging does not establish a uniform bound on unvisited
contexts. A supplied engineering bound or an independently acquired
uniform certificate retains its own full acquisition cost.

For each trial choose a coordinate uniformly, install its declared
neighborhood in the target context on a faithful local copy, and make
four successive pin--release first-exit observations from the same
requested starting bit. In the ideal Markov model the four waits at
coordinate $i$ are independent $\operatorname{Exp}(\lambda_i^u)$
conditional on the requested coordinate. Repeated physical cells need
only the uniform conditional implementation guarantees below. No
independence between imperfect trials or cells is assumed.

The context setting, four starting-state preparations, settling,
operating-field reproduction, detector operation and clock service are
all paid. In a bounded-neighborhood model this sets only the declared
local neighborhood; it does not prepare the entire global target state.
A local pin is not a free natural visit to that state. Structural closure
and faithful local replication remain premises.

\subsection{The classical inverse-gamma statistic and finite modifications}
Let $T_1,\ldots,T_4$ have rate $\lambda$ and set $S=\sum_jT_j$.
The classical unbiased rate statistic is $X=3/S$. Direct integration of
the gamma density gives
\begin{equation}
 f_S(s)=\frac{\lambda^4s^3e^{-\lambda s}}6,
 \qquad \RAE X=\lambda,
 \qquad \RAE X^2=\frac32\lambda^2.
 \label{ra:gamma}
\end{equation}
This identity is standard exponential-family estimation, not an invented
estimator; it is derived here to expose the exact variance and finite
device requirements. In particular, pooling heterogeneous waits and
inverting their overall mean would estimate a harmonic-rate quantity,
not the required arithmetic mean. The coordinate is held fixed within
each group of four, then the resulting rate statistics are averaged.

Fix $0<\eta\le1/4$. Choose a rational $a$ with
\begin{equation}
 \frac{\eta^{1/3}}{2\lambda_+}<a\le
 \frac{\eta^{1/3}}{\lambda_+},\qquad B=3/a,
 \qquad D=\frac{\lceil\log_2(192/\eta)\rceil}{\lambda_-}.
 \label{ra:clipcap}
\end{equation}
A dyadic $a\lambda_+$ satisfying these inequalities is obtained by
integer comparisons of cubes. Let
\[
 X_a=\min(3/S,B)=3/\max(S,a),\qquad
 Z=X_a\RAone_{\{\max_jT_j\le D\}}.
\]
Run the waits sequentially, stopping and returning zero as soon as one
wait is censored. The ideal output lies in $[0,B]$. The two lost-mean
bounds are
\begin{align}
 0\le\lambda-\RAE X_a
 &\le\frac{\lambda^4a^3}{24}
 \le\frac{\eta\lambda}{24},\label{ra:clipbias}\\
 0\le\RAE X_a-\RAE Z
 &\le\frac{12e^{-\lambda D}}D
 \le\frac{\eta\lambda_-}{16}.
 \label{ra:censorbias}
\end{align}
For \eqref{ra:clipbias}, integrate
$(3/s-3/a)f_S(s)$ on $[0,a]$ and bound $e^{-\lambda s}$ by one.
For \eqref{ra:censorbias}, on any $T_j>D$ one has $X_a\le3/D$;
take a union bound over four waits. Also $D\lambda_-\ge1$ and
$e^{-\lambda_-D}\le2^{-\lambda_-D}\le\eta/192$.
The clipped, censored second moment is at most $3\lambda^2/2$.

\subsection{Actual deadlines, first-event timestamps and finite arithmetic}
An actual timer must supply the following service, including start,
release, detector latency and deadline response in the error budget:
outside a declared bad-service event, an uncensored first-exit timestamp
differs from the true wait by at most $\Delta$, and each actual censor
deadline differs from $D$ by at most $\Delta$. Take
\begin{equation}
 0<\Delta\le\frac{\eta\lambda_-a^2}{1024},\qquad
 0<\omega\le\frac{\eta\lambda_-}{128},\qquad
 0\le\beta\le\frac{\eta\lambda_-}{128B}.
 \label{ra:hardware}
\end{equation}
$\omega$ bounds final numerical rounding of the rate statistic. The
total conditional TV or bad-service probability allowance $\beta$ per
trial covers finite coordinate selection, context and pin failures,
missed exits, nonideal local conditional laws and other declared
stochastic readout errors. It is a guarantee conditional on the entire
past, on the common calibration good event. Deterministic timer errors
are accounted for separately and are not represented as a false TV
approximation of a continuous timestamp by a grid.

The reported statistic is
$Y=3/\max(a,\sum_j\widehat T_j)$ if all four observations are accepted,
and zero on a censor or flagged failure, rounded within $\omega$ and
clipped to $[0,B]$. Round each timestamp with its physical errors included
in $\Delta$. Exact rational division is also allowed. Every operation
has a finite cap, including failed reads and state preparations.

The censor decision can disagree with its ideal value only when at
least one true wait lies in $[D-\Delta,D+\Delta]$, or on the bad-service
event. The ideal censor-boundary band has probability at most
$8\lambda_+\Delta$. Away from that band, the reciprocal clipping map is
$3/a^2$-Lipschitz, so its time-induced error is at most $12\Delta/a^2$.
Consequently
\begin{equation}
 |\RAE(Y\mid\mathcal F)-\RAE Z|
 \le\frac{12\Delta}{a^2}+8B\lambda_+\Delta+B\beta+\omega.
 \label{ra:timerbias}
\end{equation}
Here $\mathcal F$ contains every previous trial and acquisition outcome;
the ideal expectation includes a fresh uniform coordinate.
The band term is essential: a Lipschitz bound on a recorded sum alone
would not account for censor decisions changed by timer error.

\begin{lemma}[Uniform conditional moments of the physical trial]
On the simultaneous calibration good event, the construction satisfies
\begin{equation}
 |\RAE(Y\mid\mathcal F)-\mu_u|\le\frac{\eta\mu_u}{6},
 \qquad \RAE(Y^2\mid\mathcal F)\le4\lambda_+\mu_u
                                      \le4\lambda_+^2.
 \label{ra:trialmoments}
\end{equation}
\end{lemma}
\begin{proof}
Average \eqref{ra:clipbias} over the uniform coordinate. The clipping
loss is at most $\eta\mu_u/24$. The censor loss is at most
$\eta\lambda_-/16$. Since $\lambda_+a\le1$, the remaining errors in
\eqref{ra:timerbias} sum to at most
\[
 \eta\lambda_-\left(\frac{12+24}{1024}
                         +\frac1{128}+\frac1{128}\right)
 =\frac{13\eta\lambda_-}{256}.
\]
Thus the total is at most
$\eta\mu_u(1/24+1/16+13/256)=119\eta\mu_u/768
<\eta\mu_u/6$.

For the second moment couple to the ideal trial. Away from the censor
band and bad-service event, the output differs from $Z$ by at most
$e=12\Delta/a^2+\omega\le5\eta\lambda_-/256$.
Hence
\[
 \RAE(Y^2\mid\mathcal F)
 \le2\RAE Z^2+2e^2+B^2(8\lambda_+\Delta+\beta).
\]
The first term, averaged over the requested coordinate, is at most
$3N^{-1}\sum_i\lambda_i^2\le3\lambda_+\mu_u$.
The band contribution is at
most $72\eta\lambda_-\lambda_+/1024$.
The lower bound on $a$ gives $B<6\lambda_+\eta^{-1/3}$, so
$B^2\beta\le6\eta^{2/3}\lambda_-\lambda_+/128$.
Together with $\eta\le1/4$ these terms and $2e^2$ are strictly below
$\lambda_+\lambda_-\le\lambda_+\mu_u$.
This proves the bound without replacing the variance
by the much larger worst-case range $B^2$.
\end{proof}

The timer tolerance scales as
$\Delta=\Theta(\eta^{5/3}\lambda_-/\lambda_+^2)$ for a comparable
choice in \eqref{ra:hardware}. Its response and calibration price remain
physical costs. The number of timestamp and arithmetic bits is
$O(\log(1/\eta)+\log\kappa+\log\log(1/\eta))$ beyond the rate and
hardware scale representation. This bit count does not establish that
a device can service that resolution cheaply.

\subsection{A finite confidence bound with dependent repeated cells}
\begin{theorem}[Direct acquired arithmetic-rate certificate]
For each of $K$ declared menu settings, divide the trials into $q$ blocks
of size $m$, where
\begin{equation}
 m=\left\lceil256\kappa\eta^{-2}\right\rceil,
 \qquad
 q=\text{the smallest positive odd integer with }2^{-q}\le\delta/K.
 \label{ra:mom}
\end{equation}
Return the median $\widehat\mu_u$ of the $q$ block averages. Conditional
on the common calibration good event, with probability at least
$1-\delta$, simultaneously for all menu settings,
\begin{equation}
 |\widehat\mu_u-\mu_u|\le\eta\mu_u.
 \label{ra:certificate}
\end{equation}
In fact the proof gives the stronger constant $2\eta/3$ on that event.
The number of trials per menu is
$O(\kappa\eta^{-2}\log(K/\delta))$, independent of $N$ except for
coordinate selection, context access and their actual service costs.
No kinetic rate table, rare stationary count or global trajectory is
acquired by this assay.
\end{theorem}
\begin{proof}
For one block, write $Y_j=m_j+D_j$ with
$m_j=\RAE(Y_j\mid\mathcal F_{j-1})$ and $(D_j)$ martingale
differences. The conditional second-moment bound gives
$\RAE[(\sum_jD_j)^2\mid\mathcal F_{\rm start}]\le4m\lambda_+\mu_u$.
Conditional Chebyshev therefore implies
\[
 \RAP\left(\left|m^{-1}\sum_jD_j\right|>
        \frac{\eta\mu_u}{2}\,\middle|\,\mathcal F_{\rm start}\right)
 \le\frac{16\lambda_+}{m\eta^2\mu_u}
 \le\frac{16\kappa}{m\eta^2}\le\frac1{16}.
\]
Every $m_j$ differs from $\mu_u$ by at most $\eta\mu_u/6$.
Thus each block average is within $2\eta\mu_u/3$ except on an event
whose conditional probability, given all preceding blocks, is at most
$1/16$. For any specified set of $s$ bad blocks, repeated conditioning
bounds their joint probability by $16^{-s}$. A bad median needs at
least $\lceil q/2\rceil$ bad blocks, so its probability is at most
$2^q16^{-\lceil q/2\rceil}\le2^{-q}$. No independence of blocks is
used. Union-bound over the $K$ menus; adaptive scheduling is allowed
provided the menus and trial guarantees are the declared ones.
\end{proof}

\begin{theorem}[A smaller direct-mean schedule]
The same implemented trials can instead be averaged directly. With
\begin{equation}
 n=\left\lceil40\kappa\eta^{-2}\log(2K/\delta)\right\rceil
 \label{ra:bernsteinN}
\end{equation}
trials per menu, their ordinary average satisfies
\eqref{ra:certificate} simultaneously with probability at least
$1-\delta$, conditional on calibration. The bound again holds for
arbitrarily dependent repeated cells satisfying the conditional service
guarantees. No empirical or independence approximation is used.
\end{theorem}
\begin{proof}
For the martingale differences $D_j=Y_j-\RAE(Y_j\mid\mathcal F_{j-1})$,
$|D_j|\le B$ and the conditional variance is at most
$v=4\lambda_+\mu_u$. For $|s|B<3$, expanding the exponential and using
$k!\ge2\cdot3^{k-2}$ for $k\ge2$ gives
\[
 \RAE(e^{sD_j}\mid\mathcal F_{j-1})
 \le\exp\!\left(\frac{s^2v}{2(1-|s|B/3)}\right).
\]
Iterate this conditional bound and apply Chernoff on each side. The
resulting standard martingale Bernstein inequality is
\[
 \RAP\left(\left|n^{-1}\sum_jD_j\right|>t\right)
 \le2\exp\left[-\frac{nt^2}{2(v+Bt/3)}\right].
\]
With $t=\eta\mu_u/2$ and
$B\eta<6\lambda_+\eta^{2/3}\le6\lambda_+$, the exponent is at least
$n\eta^2\mu_u/(40\lambda_+)\ge n\eta^2/(40\kappa)$.
The conditional-mean bias adds at most $\eta\mu_u/6$, giving total
error at most $2\eta\mu_u/3$. Union-bound over the menus.
\end{proof}

If common apparatus or structural calibration fails with probability
$\delta_{\rm cal}$, the unconditional confidence is at least
$1-\delta_{\rm cal}-\delta$. This is distinct from a per-trial conditional
error; reusing a persistently miscalibrated device does not generate new
independent confidence. A common drift or hidden cell memory not covered
by \eqref{ra:trialmoments} invalidates the stated martingale guarantee.

The reported interval for $\mu_u$ is
$[\widehat\mu_u/(1+\eta),\widehat\mu_u/(1-\eta)]$, with outward
rounding; multiply by $N$ for $\Lambda_u$. Choosing the smallest
$\widehat\mu_u$ yields
\[
 \mu_{\widehat u}\le\frac{1+\eta}{1-\eta}\min_u\mu_u.
\]
If the resulting intervals are disjoint, their ordering is certified.
The separate mathematical bridge from boundary escape to the intended
global first-event objective must still be proved at its stated scope;
the local assay alone does not equate stationary flux with lifetime.

The direct-mean implementation is streaming: it stores four current
timestamps, a sum and count, a coordinate index, and the $K$ final
intervals. The median variant additionally stores its $q$ block means.
The arithmetic bit lengths are polynomial in rate and hardware bit
length, $\log N$, $\log n$, and requested precision. Known coordinate
labels and neighborhood access, apparatus cells, context memory and
addressing latency remain priced resources; no free global wiring or
new physical copies are inferred from this small statistical memory.

\subsection{Full dwell and reset bill, including discarded trials}
Let $s=e^{-\lambda D}$ and $q_\lambda=1-s$ for a requested rate.
The ideal sequentially censored trial has expected active dwell
\begin{equation}
 \RAE W_\lambda
 =\frac{q_\lambda}{\lambda}\sum_{j=0}^3q_\lambda^j
 \le\frac4\lambda,
 \qquad W_\lambda\le4D.
 \label{ra:dwell}
\end{equation}
The expected number of pin--release observations is
$\sum_{j=0}^3q_\lambda^j\le4$.
The probability of a discarded censored trial is $1-q_\lambda^4$.
Its entire wasted dwell, including the successful waits before its first
censor, has the exact expectation
\begin{equation}
 \RAE[W_\lambda\RAone_{\rm censored}]
 =\sum_{j=1}^4q_\lambda^{j-1}s
 \left[D+(j-1)\left(\frac1\lambda-
                  \frac{Ds}{q_\lambda}\right)\right]
 \le16D e^{-\lambda D}.
 \label{ra:waste}
\end{equation}
The formulas follow by conditioning on the index of the first censored
wait; $\RAE[T\mid T\le D]=\lambda^{-1}-Ds/q_\lambda$.
Here $q_\lambda>0$ because $D>0$.

With the actual timestamp/deadline service, each normal trial has active
dwell at most $4(D+\Delta)$ and expected active dwell at most
$4/\lambda_-+4\Delta$. A conservative addition
$4\beta(D+\Delta)$ accounts for conditional bad-service branches, whose
timeouts are also hard-capped. Discarded active dwell is bounded by
$16(D+\Delta)e^{-\lambda_-(D-\Delta)}+4\beta(D+\Delta)$ per trial.
All pin, clamp, setup, settling and readout times are additional. These
bounds include unsuccessful exposure; censoring does not make it free.

With $n=mq$ trials for the median version, or \eqref{ra:bernsteinN}
trials for the direct mean, the worst-case active dwell is
$4n(D+\Delta)$ and the expected active dwell is at most
$n[4/\lambda_-+4\Delta+4\beta(D+\Delta)]$.
The complete serial wall-time bound adds, for each trial,
\[
 C_{\rm coordinate}+C_{\rm context}+C_{\rm operating\ setting}
 +4(C_{\rm pin}+C_{\rm settle}+C_{\rm detector}+C_{\rm timer})
 +C_{\rm arithmetic},
\]
plus common calibration, inter-menu finite slew and final deployment.
Use the actual parallel service architecture to convert total work and
component exposure to elapsed time; parallel copies are additional paid
resources. No unpriced simultaneity is assumed. The existing safe mission
preparation and deployment bill are separate from these local target-context
assays and are not removed by a cheaper decision statistic.

The apparatus must detect the \emph{first exit}, not merely read a late
endpoint after possible returns. A failed event record may be flagged
and charged in $\beta$; an unbounded missed-event mechanism cannot be
silently reinterpreted as an exponential wait. If the requested exit
floor is small, its numerical $\lambda_-^{-1}$ exposure is retained.

\section{A matched decision-exposure lower bound at the same local access}
\label{sec:rate_lower}
The quadratic accuracy cost cannot generally be removed by observing more
of a local trajectory, choosing fields adaptively, or granting free
parameter-independent pins. The following argument is an elementary
continuous-time likelihood comparison, in the tradition of sequential
change-of-measure lower bounds. It does not claim that inverse-gamma
estimation or sequential testing is new.

\begin{theorem}[Two-menu decision requires quadratic local exposure]
Fix $0<\eta\le1/4$, $0<\delta<1/2$, $\lambda_0>0$, any $N\ge1$,
and a known $q_*\in(0,1/4]$.
There are two models of $N$ independent bits and two operating menus
such that:
\begin{enumerate}
\item both menus in both models have the same product invariant law
with one-bit probability $q_*$, and the singleton target $b=1^N$ has
$p=q_*^N$;
\item every local \emph{boundary recovery} rate $Q(b,b^{(i)})$ is in
$[\lambda_0(1-2\eta),\lambda_0(1+2\eta)]$;
\item the two models reverse which menu minimizes the mean first-event
time, and also reverse which maximizes it, with a multiplicative gap
larger than $1+\eta$;
\item any adaptive decision procedure with error at most $\delta$ on
both models must have expected total active local exposure
\begin{equation}
 \RAE T_{\rm exposure}\ge
 \frac{(1-2\delta)\log((1-\delta)/\delta)}
      {64\lambda_0\eta^2}
 \label{ra:lower}
\end{equation}
under each model.
\end{enumerate}
The procedure may observe complete continuous local jump paths, use
arbitrary predictable calibrated fields, and select operating menus and
starting bits adaptively. Free resets may reveal their commanded state
but no additional parameter information.
With parallel cells $T_{\rm exposure}$ is the sum of their active
durations, not wall time. The permitted parameter-dependent observations
are the local state/jump records; no extra rate-revealing actuator
side-channel is supplied.
\end{theorem}
\begin{proof}
Put $\theta=2\eta\le1/2$ and
$\lambda_\pm=\lambda_0(1\pm\theta)$.
In model zero, menu one has single-bit $1\to0$ recovery rate $\lambda_+$
and menu two has recovery rate $\lambda_-$. The corresponding $0\to1$
failure rate is $\lambda_\pm q_* /(1-q_*)$; for $q_*=1/4$ it is
$\lambda_\pm/3$. Swap the menus in model one.
Both global generators are common scalar multiples of the same
independent-bit generator. Their equilibrium-safe means therefore scale
exactly inversely with their local flip rate. The ratio is
$(1+\theta)/(1-\theta)>1+\eta$ for every $N$, and the invariant target
mass remains $q_*^N$. Given a strict reliable-family bound
$\bar q<1/2$, choosing $0<q_*\le\min(\bar q/2,1/4)$ places this hard
pair strictly within that family.

To include arbitrary experimental fields, let a field $z$ change the
known conditional probability to $\sigma(\operatorname{logit}(q_*)+z)$,
while the attempt clock in that menu remains $\lambda_\pm/(1-q_*)$.
This is a valid subclass of the
existing access model. The two possible jump intensities at any history
have the form $a=\lambda_0(1+\theta)v/(1-q_*)$ and
$b=\lambda_0(1-\theta)v/(1-q_*)$ in one order, with the same known
$v\in[0,1]$ determined by the field and current bit.
For nonnegative intensities,
\[
 a\log(a/b)-a+b\le (a-b)^2/b
 \le\frac{4\lambda_0\theta^2}{(1-q_*)(1-\theta)}
 \le16\lambda_0\theta^2=64\lambda_0\eta^2.
\]
The same bound applies in the opposite direction; the zero-intensity
limit is interpreted continuously. The controlled point-process
likelihood chain rule integrates this divergence rate over the total
active local exposure. Predictable choices and parameter-independent
resets use the same conditional kernels under the two models and add
no divergence. Consequently
\[
 \RAKL(\mathsf P_0\Vert\mathsf P_1)
 \le64\lambda_0\eta^2\RAE_0T_{\rm exposure},
\]
and conversely with the model indices exchanged. Bounded exposure gives
this directly. For integrable adaptive stopping, stop first at a finite
exposure cap and pass to the limit in the likelihood compensator.

A correct $(1+\eta)$ menu decision must select the appropriate menu in
each model, so it is a binary test with both errors at most $\delta$.
Data processing to that decision gives
\[
 \RAKL(\mathsf P_0\Vert\mathsf P_1)
 \ge \operatorname{kl}(1-\delta,\delta)
 =(1-2\delta)\log((1-\delta)/\delta).
\]
Combine the inequalities in both directions.
\end{proof}

At fixed contrast, the positive assay's expected active exposure is
$O(\lambda_-^{-1}\eta^{-2}\log(K/\delta))$, matching the lower bound's
accuracy and confidence dependence for two menus up to constants and
the stated contrast. The hard worst-case censor cap adds
$\log(1/\eta)$; pins, service latencies and calibration are additional
resources and are not lower-bounded by \eqref{ra:lower}.
The lower bound is uniform in the exponentially small global target
probability. It does not assert optimal dependence on contrast, menu
count, detector price or physical clock resolution.

\section{Ordinary methods, the new interface and retained limits}
The statistic $(n-1)/\sum_{j=1}^nT_j$ for an exponential rate is
classical; the gamma integral above proves the needed $n=4$ instance.
Likelihood-based exponential-rate estimation is an appropriate strong
ordinary comparator. A method that estimates every local rate can also
reach quadratic precision cost per rate, but needlessly acquires $N$
parameters when the decision requires only their average. Conversely,
an ordinary method that groups repeated waits within a randomly chosen
context and then averages the classical unbiased statistics implements
the same direct representation and can match this result. No claim is
made that every conventional estimator has cubic cost.

In particular, a small-time endpoint estimator that divides an observed
exit probability by its hold time uses a time of order $\eta$ to suppress
linearization bias. Its reset/read count can scale as $\eta^{-3}$ even
though its total exposure already scales quadratically. The present
grouped first-exit statistic removes that particular extra reset factor,
retains a finite physical clock construction, and estimates the
arithmetic rate average without a rate table. It does not improve on a
stronger grouped-wait method already using this identity.

The new contribution for the chosen decision is the complete interface:
requested operating contexts, an average of heterogeneous fast-side
rates, finite censor/clip/timing errors including deadline boundary bands,
dependence-safe confidence, all discarded exposure, and a matched
controlled-observation lower bound. Its use in a global rare-event
intervention still relies on the separate target-set and spectral or
potential-theoretic decision theorem. Nothing in the local data certifies
that unobserved failure mechanisms are absent.
